\documentclass[10pt,a4paper]{article}
\usepackage[margin=1.25cm]{geometry}
\usepackage{amsmath,amssymb,mathtools}
\usepackage{multicol}
\usepackage{array,booktabs}
\usepackage{xcolor}
\usepackage{enumitem}
\usepackage{titlesec}
\usepackage{fancyhdr}
\usepackage{tcolorbox}
\usepackage{lmodern}
\setlength{\parindent}{0pt}
\setlength{\parskip}{2pt}
\setlist[itemize]{leftmargin=1.2em,itemsep=1pt,topsep=1pt}
\setlist[enumerate]{leftmargin=1.4em,itemsep=1pt,topsep=1pt}
\titleformat{\section}{\large\bfseries}{}{0pt}{}
\titleformat{\subsection}{\normalsize\bfseries}{}{0pt}{}
\pagestyle{fancy}
\fancyhf{}
\lhead{Calculus III Final Exam Cheatsheet}
\rhead{MTH253}
\cfoot{\thepage}
\newtcolorbox{rulebox}{colback=gray!6,colframe=black!35,boxrule=0.5pt,arc=1mm,left=2mm,right=2mm,top=1mm,bottom=1mm}
\newcommand{\boxeq}[1]{\[\boxed{#1}\]}

\begin{document}

\begin{center}
{\LARGE\bfseries Calculus III Final Exam Cheatsheet}\\[2mm]
{\small Directional derivative - Absolute extrema/Lagrange - Multiple integrals - Polar/Spherical - Scalar line integral - Green's theorem - Flux}
\end{center}

\begin{rulebox}
\textbf{Fast recognition.}
Direction/change $\to$ directional derivative. \quad
Absolute max/min on closed $D$ $\to$ interior + boundary + corners. \quad
Constraint $g=c$ $\to$ Lagrange. \quad
Circles/disks $\to$ polar. \quad
Spheres/cones $\to$ spherical. \quad
$\int_C f\,ds$ $\to$ scalar line integral. \quad
$\oint_C P\,dx+Q\,dy$ with closed $C$ $\to$ Green. \quad
$\iint_S \mathbf F\cdot\mathbf n\,dS$ $\to$ flux.
\end{rulebox}

\begin{multicols}{2}

\section*{1. Directional Derivative}
For $f(x,y,z)$,
\boxeq{\nabla f=\langle f_x,f_y,f_z\rangle}
For a given direction vector $\mathbf v$,
\boxeq{\mathbf u=\frac{\mathbf v}{|\mathbf v|}}
Then
\boxeq{D_{\mathbf u}f(P)=\nabla f(P)\cdot\mathbf u}
If direction is from $A$ to $B$:
\boxeq{\mathbf v=B-A,\qquad \mathbf u=\frac{B-A}{|B-A|}}
If a direction angle $\theta$ is given in the plane:
\boxeq{\mathbf u=\langle\cos\theta,\sin\theta\rangle}
Fastest increase:
\boxeq{\text{direction}=\frac{\nabla f}{|\nabla f|},\qquad \text{rate}=|\nabla f|}
Fastest decrease:
\boxeq{\text{direction}=-\frac{\nabla f}{|\nabla f|},\qquad \text{rate}=-|\nabla f|}
\textbf{Rule of thumb:} if the problem gives a raw direction vector, \emph{normalize first}. If $\mathbf u\perp \nabla f$, then $D_{\mathbf u}f=0$.

\section*{2. Absolute Max/Min + Lagrange}
For a continuous $f$ on a closed, bounded region $D$, absolute extrema exist.

\textbf{Interior candidates:}
\boxeq{f_x=0,\qquad f_y=0}
Keep only solutions inside $D$.

\textbf{Boundary constraint:} if $g(x,y)=c$,
\boxeq{\nabla f=\lambda\nabla g}
that is,
\boxeq{f_x=\lambda g_x,\qquad f_y=\lambda g_y,\qquad g=c}
In 3D add $f_z=\lambda g_z$.

\textbf{Final comparison checklist:}
\begin{enumerate}
\item Interior critical points.
\item Every smooth boundary piece (Lagrange or parametrization).
\item Corners/endpoints.
\item Evaluate $f$ at all candidates.
\item Largest value = absolute max; smallest = absolute min.
\end{enumerate}
\textbf{Trap:} Lagrange handles the constrained boundary; it does \emph{not} replace checking the interior.

\section*{3. Double and Triple Integrals}
General double integral:
\boxeq{\iint_D f(x,y)\,dA}
Type I region:
\boxeq{\iint_D f\,dA=\int_a^b\int_{g_1(x)}^{g_2(x)} f(x,y)\,dy\,dx}
Area:
\boxeq{\iint_D 1\,dA=\operatorname{Area}(D)}
If $f\ge 0$, $\iint_D f\,dA$ is volume under $z=f(x,y)$.

Triple integral:
\boxeq{\iiint_E f(x,y,z)\,dV}
Volume:
\boxeq{\iiint_E1\,dV=\operatorname{Vol}(E)}
If $z_{\min}(x,y)\le z\le z_{\max}(x,y)$:
\boxeq{\iiint_E f\,dV=\iint_D\int_{z_{\min}}^{z_{\max}} f\,dz\,dA}
\textbf{Rule:} read iterated integrals from inside out.

\section*{4. Polar Coordinates}
Use for disks, circles, annuli, sectors, or $x^2+y^2$.
\boxeq{x=r\cos\theta,\quad y=r\sin\theta,\quad x^2+y^2=r^2}
\boxeq{dA=r\,dr\,d\theta}
Therefore
\boxeq{\iint_D f\,dA=\int_{\alpha}^{\beta}\int_{r_1(\theta)}^{r_2(\theta)}f(r\cos\theta,r\sin\theta)\,r\,dr\,d\theta}
Standard angles:
\begin{itemize}
\item Full circle: $0\le\theta\le2\pi$.
\item Upper half: $0\le\theta\le\pi$.
\item First quadrant: $0\le\theta\le\pi/2$.
\item Right half: $-\pi/2\le\theta\le\pi/2$.
\end{itemize}
Useful conversions:
\[
x^2+y^2=R^2\Rightarrow r=R.
\]
For other boundaries, substitute $x=r\cos\theta$ and $y=r\sin\theta$ and solve directly for $r$.
\textbf{Trap:} never forget the Jacobian factor $r$.

\section*{5. Spherical Coordinates}
Use for spheres, hemispheres, cones with vertex at the origin, radial functions.
\boxeq{x=\rho\sin\phi\cos\theta}
\boxeq{y=\rho\sin\phi\sin\theta}
\boxeq{z=\rho\cos\phi}
\boxeq{x^2+y^2+z^2=\rho^2,\qquad \sqrt{x^2+y^2}=\rho\sin\phi}
Volume element:
\boxeq{dV=\rho^2\sin\phi\,d\rho\,d\phi\,d\theta}
Meaning:
\begin{itemize}
\item $\rho$: distance from origin.
\item $\theta$: rotation around the $z$-axis.
\item $\phi$: angle measured down from the positive $z$-axis.
\end{itemize}
Whole sphere: $0\le\rho\le R$, $0\le\phi\le\pi$, $0\le\theta\le2\pi$.
Upper hemisphere: $0\le\phi\le\pi/2$.

Cone trick: for $z=\sqrt{x^2+y^2}$,
\[
\rho\cos\phi=\rho\sin\phi\Rightarrow \phi=\frac{\pi}{4}.
\]
\textbf{Memory rule:} sphere $\to\rho$ bound; cone $\to\phi$ bound; rotation $\to\theta$ bound.

\section*{6. Scalar Line Integral (Type 1)}
Parametrize the curve:
\boxeq{\mathbf r(t)=\langle x(t),y(t),z(t)\rangle}
Then
\boxeq{\mathbf r'(t)=\langle x'(t),y'(t),z'(t)\rangle}
\boxeq{ds=|\mathbf r'(t)|\,dt}
where
\boxeq{|\mathbf r'(t)|=\sqrt{x'(t)^2+y'(t)^2+z'(t)^2}}
Core formula:
\boxeq{\int_C f\,ds=\int_a^b f(\mathbf r(t))|\mathbf r'(t)|\,dt}
In 2D simply omit $z'(t)^2$.

Special case:
\boxeq{\int_C1\,ds=\text{length of }C}
If $f$ is linear density, $\int_C f\,ds$ is mass of a wire.

\textbf{Orientation:}
\boxeq{\int_C f\,ds=\int_{-C}f\,ds}
Type 1 does not care which way the curve is traversed.

Useful parametrizations:
\[
\mathbf r(t)=A+t(B-A),\ 0\le t\le1
\]
for a line segment, and
\[
x=a+R\cos t,\qquad y=b+R\sin t
\]
for a circle.

\section*{7. Green's Theorem}
For $\mathbf F=(P,Q)$ and positively oriented (counterclockwise) closed curve $C=\partial D$:
\boxeq{\oint_C P\,dx+Q\,dy=\iint_D\left(\frac{\partial Q}{\partial x}-\frac{\partial P}{\partial y}\right)\,dA}
Memorize:
\boxeq{Q_x-P_y}
Positive orientation = counterclockwise = region stays on your left.

If $C$ is clockwise:
\boxeq{\oint_C P\,dx+Q\,dy=-\iint_D(Q_x-P_y)\,dA}
\textbf{Workflow:} identify $P,Q$ $\to$ compute $Q_x-P_y$ $\to$ find region $D$ $\to$ integrate.

\textbf{Trap:} Green applies to a \emph{vector/type-2 line integral}, not to $\int_C f\,ds$.

\section*{8. Surface Integral Type 2 (Flux)}
For vector field
\boxeq{\mathbf F=(P,Q,R)}
flux through an oriented surface $S$ is
\boxeq{I=\iint_S\mathbf F\cdot\mathbf n\,dS}
Lecture notation:
\boxeq{I=\iint_S P\,dy\,dz+Q\,dz\,dx+R\,dx\,dy}
$\mathbf F\cdot\mathbf n$ extracts the component of the field normal to the surface.

If a surface is implicit, $G(x,y,z)=0$:
\boxeq{\nabla G=\langle G_x,G_y,G_z\rangle\text{ is normal to }S}
Unit normal:
\boxeq{\mathbf n=\pm\frac{\nabla G}{\|\nabla G\|}}

\subsection*{Graph shortcut: $z=g(x,y)$}
For \textbf{upward} orientation:
\boxeq{\mathbf n\,dS=\langle-g_x,-g_y,1\rangle\,dx\,dy}
Hence
\boxeq{\iint_S\mathbf F\cdot\mathbf n\,dS=\iint_D(-Pg_x-Qg_y+R)\,dx\,dy}
For \textbf{downward} orientation:
\boxeq{\mathbf n\,dS=\langle g_x,g_y,-1\rangle\,dx\,dy}
Hence
\boxeq{\iint_S\mathbf F\cdot\mathbf n\,dS=\iint_D(Pg_x+Qg_y-R)\,dx\,dy}
Evaluate $P,Q,R$ on the surface $z=g(x,y)$.

\textbf{Why no normalization?} In the product $\mathbf n\,dS$, the unit-normal denominator $\sqrt{1+g_x^2+g_y^2}$ cancels the same factor in $dS$.

\subsection*{Parametric surface}
If $\mathbf r(u,v)$ parametrizes $S$,
\boxeq{\iint_S\mathbf F\cdot\mathbf n\,dS=\iint_D\mathbf F(\mathbf r(u,v))\cdot(\mathbf r_u\times\mathbf r_v)\,du\,dv}
If orientation is opposite, multiply the normal by $-1$.

\section*{9. Physical / Geometric Meaning}
\begin{itemize}
\item $\int_a^b f(x)\,dx$: accumulation along an interval.
\item $\iint_D f\,dA$: accumulation over area; if $f\ge0$, volume under a surface.
\item $\iiint_E f\,dV$: accumulation over volume; $f=1$ gives volume.
\item $\int_C f\,ds$: accumulation along curve length; $f=1$ gives curve length.
\item $\iint_S \mathbf F\cdot\mathbf n\,dS$: net normal flow (flux) through a surface.
\end{itemize}
Flux signs (for chosen normal):
\begin{itemize}
\item $\Phi>0$: net flow in the normal direction.
\item $\Phi<0$: net flow opposite the normal direction.
\item $\Phi=0$: tangent flow or cancellation of positive/negative contributions; field need not be zero.
\end{itemize}

\section*{10. Last-Minute Traps}
\begin{enumerate}
\item Directional derivative: normalize the direction vector.
\item Absolute extrema: check interior + boundary + corners.
\item Polar: $dA=r\,dr\,d\theta$.
\item Spherical: $dV=\rho^2\sin\phi\,d\rho\,d\phi\,d\theta$.
\item $\theta$ rotates around $z$; $\phi$ is measured down from $+z$.
\item Scalar line integral: include $ds=|\mathbf r'(t)|dt$.
\item Green: curve must be closed; CCW is positive.
\item Flux: orientation matters; reversing the normal changes the sign.
\item Graph-flux shortcut already combines $\mathbf n$ and $dS$; do not normalize again.
\end{enumerate}

\end{multicols}

\begin{rulebox}
\textbf{10 formulas to memorize first:}\quad
$\nabla f$, \quad
$\mathbf u=\mathbf v/|\mathbf v|$, \quad
$D_{\mathbf u}f=\nabla f\cdot\mathbf u$, \quad
$\nabla f=\lambda\nabla g$, \quad
$dA=r\,dr\,d\theta$, \quad
$dV=\rho^2\sin\phi\,d\rho\,d\phi\,d\theta$, \quad
$\int_C f\,ds=\int f(\mathbf r(t))|\mathbf r'(t)|dt$, \quad
$\oint_C Pdx+Qdy=\iint_D(Q_x-P_y)dA$, \quad
$I=\iint_S\mathbf F\cdot\mathbf n\,dS$, \quad
$z=g(x,y):\ \text{up }(-g_x,-g_y,1),\ \text{down }(g_x,g_y,-1)$.
\end{rulebox}

\end{document}
